
\documentclass{article}
\usepackage{colortbl}
\usepackage{makecell}
\usepackage{multirow}
\usepackage{supertabular}

\begin{document}

\newcounter{utterance}

\centering \large Interaction Transcript for game `dond', experiment `coop\_en', episode 0 with baseline\_clp\_chat.
\vspace{24pt}

{ \footnotesize  \setcounter{utterance}{1}
\setlength{\tabcolsep}{0pt}
\begin{supertabular}{c@{$\;$}|p{.15\linewidth}@{}p{.15\linewidth}p{.15\linewidth}p{.15\linewidth}p{.15\linewidth}p{.15\linewidth}}
    \# & $\;$A & \multicolumn{4}{c}{Game Master} & $\;\:$B\\
    \hline

    \theutterance \stepcounter{utterance}  
    & & \multicolumn{4}{p{0.6\linewidth}}{
        \cellcolor[rgb]{0.9,0.9,0.9}{
            \makecell[{{p{\linewidth}}}]{
                \texttt{\tiny{[P1$\langle$GM]}}
                \texttt{You are playing a cooperative negotiation game in which you have to agree on how to divide a set of items among you and another player.} \\
\\ 
\texttt{Rules:} \\
\texttt{(a) You and the other player are given a set of items. Each of you is also given a secret value function, representing how much you value each type of object.} \\
\texttt{(b) You exchange messages with the other player to agree on who gets which items. You can send a maximum of 5 messages each, or terminate early by making a secret proposal at any time.} \\
\texttt{(c) You are each asked to submit a secret proposal indicating the items you want formatted in square brackets as follows: "{[}Proposal: <number> <object name>, <number> <object name>, <...>{]}"} \\
\texttt{(d) If your proposals are complementary, i.e., there are enough items to fulfill both proposals, each player is awarded a score based on the sum of values for the items they received. Otherwise, both of you get zero points.} \\
\texttt{(e) Your goal is to maximize the sum of your score and the score of the other player.} \\
\\ 
\texttt{Let us start.} \\
\\ 
\texttt{The set of available items is:} \\
\\ 
\texttt{1 picture, 1 fork, 1 pan, 1 shoe, 1 sponge} \\
\\ 
\texttt{Your secret value function is:} \\
\\ 
\texttt{picture: 5, fork: 1, pan: 0, shoe: 3, sponge: 1} \\
\\ 
\texttt{IMPORTANT: Your messages, unless it is the secret proposal, are directly transmitted to the other player, so do not include any response to the rules or text announcing your message. To make a secret proposal, use the indicated format. Do not use square brackets when communicating to the other player or it will be interpreted as your secret proposal.} \\
            }
        }
    }
    & \\ \\

    \theutterance \stepcounter{utterance}  
    & \multicolumn{4}{p{0.6\linewidth}}{
        \cellcolor[rgb]{0.9,0.9,0.9}{
            \makecell[{{p{\linewidth}}}]{
                \texttt{\tiny{[P1$\rangle$GM]}}
                \texttt{I value the picture most (5), followed by the shoe (3), then the sponge (1), and the fork (1). The pan is worth 0 to me, so I don't care about it. I want to maximize the total value, so I should aim for a division where both of us get items we value highly. I'll start by sending a message to the other player to see what they value.} \\
\\ 
\texttt{I value the picture the most, followed by the shoe, then the sponge and fork. I don't care about the pan. What do you value most?} \\
            }
        }
    }
    & & \\ \\

    \theutterance \stepcounter{utterance}  
    & & & \multicolumn{2}{p{0.3\linewidth}}{
        \cellcolor[rgb]{0.9,0.9,0.9}{
            \makecell[{{p{\linewidth}}}]{
                \texttt{\tiny{[GM$|$GM]}}
                \texttt{continue} \\
            }
        }
    }
    & & \\ \\

    \theutterance \stepcounter{utterance}  
    & & & \multicolumn{2}{p{0.3\linewidth}}{
        \cellcolor[rgb]{0.9,0.9,0.9}{
            \makecell[{{p{\linewidth}}}]{
                \texttt{\tiny{[GM$|$GM]}}
                \texttt{I value the picture most (5), followed by the shoe (3), then the sponge (1), and the fork (1). The pan is worth 0 to me, so I don't care about it. I want to maximize the total value, so I should aim for a division where both of us get items we value highly. I'll start by sending a message to the other player to see what they value.} \\
\\ 
\texttt{I value the picture the most, followed by the shoe, then the sponge and fork. I don't care about the pan. What do you value most?} \\
            }
        }
    }
    & & \\ \\

    \theutterance \stepcounter{utterance}  
    & & \multicolumn{4}{p{0.6\linewidth}}{
        \cellcolor[rgb]{0.9,0.9,0.9}{
            \makecell[{{p{\linewidth}}}]{
                \texttt{\tiny{[GM$\rangle$P2]}}
                \texttt{You are playing a cooperative negotiation game in which you have to agree on how to divide a set of items among you and another player.} \\
\\ 
\texttt{Rules:} \\
\texttt{(a) You and the other player are given a set of items. Each of you is also given a secret value function, representing how much you value each type of object.} \\
\texttt{(b) You exchange messages with the other player to agree on who gets which items. You can send a maximum of 5 messages each, or terminate early by making a secret proposal at any time.} \\
\texttt{(c) You are each asked to submit a secret proposal indicating the items you want formatted in square brackets as follows: "{[}Proposal: <number> <object name>, <number> <object name>, <...>{]}"} \\
\texttt{(d) If your proposals are complementary, i.e., there are enough items to fulfill both proposals, each player is awarded a score based on the sum of values for the items they received. Otherwise, both of you get zero points.} \\
\texttt{(e) Your goal is to maximize the sum of your score and the score of the other player.} \\
\\ 
\texttt{Let us start.} \\
\\ 
\texttt{The set of available items is:} \\
\\ 
\texttt{1 picture, 1 fork, 1 pan, 1 shoe, 1 sponge} \\
\\ 
\texttt{Your secret value function is:} \\
\\ 
\texttt{picture: 0, fork: 4, pan: 2, shoe: 0, sponge: 4} \\
\\ 
\texttt{IMPORTANT: Your messages, unless it is the secret proposal, are directly transmitted to the other player, so do not include any response to the rules or text announcing your message. To make a secret proposal, use the indicated format. Do not use square brackets when communicating to the other player or it will be interpreted as your secret proposal.} \\
\\ 
\\ 
\texttt{I value the picture most (5), followed by the shoe (3), then the sponge (1), and the fork (1). The pan is worth 0 to me, so I don't care about it. I want to maximize the total value, so I should aim for a division where both of us get items we value highly. I'll start by sending a message to the other player to see what they value.} \\
\\ 
\texttt{I value the picture the most, followed by the shoe, then the sponge and fork. I don't care about the pan. What do you value most?} \\
            }
        }
    }
    & \\ \\

    \theutterance \stepcounter{utterance}  
    & & & \multicolumn{4}{p{0.6\linewidth}}{
        \cellcolor[rgb]{0.9,0.9,0.9}{
            \makecell[{{p{\linewidth}}}]{
                \texttt{\tiny{[GM$\langle$P2]}}
                \texttt{I value the fork the most (4), followed by the pan (2), then the sponge (1), and the picture (0). I don't care about the shoe. I want to maximize the total value, so I should aim for a division where both of us get items we value highly. I'll start by sending a message to the other player to see what they value.} \\
\\ 
\texttt{I value the fork most, then the pan, then the sponge, and the picture and shoe are worth 0 to me. What do you value most?} \\
            }
        }
    }
     \\ \\

    \theutterance \stepcounter{utterance}  
    & & & \multicolumn{2}{p{0.3\linewidth}}{
        \cellcolor[rgb]{0.9,0.9,0.9}{
            \makecell[{{p{\linewidth}}}]{
                \texttt{\tiny{[GM$|$GM]}}
                \texttt{continue} \\
            }
        }
    }
    & & \\ \\

    \theutterance \stepcounter{utterance}  
    & & & \multicolumn{2}{p{0.3\linewidth}}{
        \cellcolor[rgb]{0.9,0.9,0.9}{
            \makecell[{{p{\linewidth}}}]{
                \texttt{\tiny{[GM$|$GM]}}
                \texttt{I value the fork the most (4), followed by the pan (2), then the sponge (1), and the picture (0). I don't care about the shoe. I want to maximize the total value, so I should aim for a division where both of us get items we value highly. I'll start by sending a message to the other player to see what they value.} \\
\\ 
\texttt{I value the fork most, then the pan, then the sponge, and the picture and shoe are worth 0 to me. What do you value most?} \\
            }
        }
    }
    & & \\ \\

    \theutterance \stepcounter{utterance}  
    & & \multicolumn{4}{p{0.6\linewidth}}{
        \cellcolor[rgb]{0.9,0.9,0.9}{
            \makecell[{{p{\linewidth}}}]{
                \texttt{\tiny{[P1$\langle$GM]}}
                \texttt{I value the fork the most (4), followed by the pan (2), then the sponge (1), and the picture (0). I don't care about the shoe. I want to maximize the total value, so I should aim for a division where both of us get items we value highly. I'll start by sending a message to the other player to see what they value.} \\
\\ 
\texttt{I value the fork most, then the pan, then the sponge, and the picture and shoe are worth 0 to me. What do you value most?} \\
            }
        }
    }
    & \\ \\

    \theutterance \stepcounter{utterance}  
    & \multicolumn{4}{p{0.6\linewidth}}{
        \cellcolor[rgb]{0.9,0.9,0.9}{
            \makecell[{{p{\linewidth}}}]{
                \texttt{\tiny{[P1$\rangle$GM]}}
                \texttt{I value the picture the most (5), followed by the shoe (3), then the sponge (1), and the fork (1). The pan is worth 0 to me, so I don't care about it. I want to maximize the total value, so I should aim for a division where both of us get items we value highly. I'll start by sending a message to the other player to see what they value.} \\
\\ 
\texttt{I value the picture the most, followed by the shoe, then the sponge and fork. I don't care about the pan. What do you value most?} \\
            }
        }
    }
    & & \\ \\

    \theutterance \stepcounter{utterance}  
    & & & \multicolumn{2}{p{0.3\linewidth}}{
        \cellcolor[rgb]{0.9,0.9,0.9}{
            \makecell[{{p{\linewidth}}}]{
                \texttt{\tiny{[GM$|$GM]}}
                \texttt{continue} \\
            }
        }
    }
    & & \\ \\

    \theutterance \stepcounter{utterance}  
    & & & \multicolumn{2}{p{0.3\linewidth}}{
        \cellcolor[rgb]{0.9,0.9,0.9}{
            \makecell[{{p{\linewidth}}}]{
                \texttt{\tiny{[GM$|$GM]}}
                \texttt{I value the picture the most (5), followed by the shoe (3), then the sponge (1), and the fork (1). The pan is worth 0 to me, so I don't care about it. I want to maximize the total value, so I should aim for a division where both of us get items we value highly. I'll start by sending a message to the other player to see what they value.} \\
\\ 
\texttt{I value the picture the most, followed by the shoe, then the sponge and fork. I don't care about the pan. What do you value most?} \\
            }
        }
    }
    & & \\ \\

    \theutterance \stepcounter{utterance}  
    & & \multicolumn{4}{p{0.6\linewidth}}{
        \cellcolor[rgb]{0.9,0.9,0.9}{
            \makecell[{{p{\linewidth}}}]{
                \texttt{\tiny{[GM$\rangle$P2]}}
                \texttt{I value the picture the most (5), followed by the shoe (3), then the sponge (1), and the fork (1). The pan is worth 0 to me, so I don't care about it. I want to maximize the total value, so I should aim for a division where both of us get items we value highly. I'll start by sending a message to the other player to see what they value.} \\
\\ 
\texttt{I value the picture the most, followed by the shoe, then the sponge and fork. I don't care about the pan. What do you value most?} \\
            }
        }
    }
    & \\ \\

    \theutterance \stepcounter{utterance}  
    & & & \multicolumn{4}{p{0.6\linewidth}}{
        \cellcolor[rgb]{0.9,0.9,0.9}{
            \makecell[{{p{\linewidth}}}]{
                \texttt{\tiny{[GM$\langle$P2]}}
                \texttt{I value the fork most (4), then the sponge (4), then the pan (2), and the picture and shoe are worth 0 to me. I don't care about the picture or shoe. I want to maximize total value, so I should aim for a division where we each get high{-}value items. Let's see if we can coordinate.} \\
\\ 
\texttt{I propose that I take the fork and the sponge (total value 8), and you take the pan (value 2). That gives us a total of 10. But I don't know if you value the pan. Alternatively, if you value the pan, maybe you should take it, and I take the fork and sponge. But I don't want to waste items.} \\
\\ 
\texttt{Wait — let me think again. You value the fork (4), pan (2), sponge (4), picture (0), shoe (0). I value fork (4), pan (2), sponge (4), picture (0), shoe (0). So both of us value fork, pan, and sponge. But we both value the fork and sponge the most.} \\
\\ 
\texttt{But we can't both have the same item. So we need to split them.} \\
\\ 
\texttt{Let’s try to maximize total value. The total value of all items is:} \\
\\ 
\texttt{{-} Fork: 4 (me) + 4 (you) = 8  } \\
\texttt{{-} Pan: 2 (me) + 2 (you) = 4  } \\
\texttt{{-} Sponge: 4 (me) + 1 (you) = 5  } \\
\texttt{{-} Picture: 0 + 5 = 5  } \\
\texttt{{-} Shoe: 0 + 3 = 3  } \\
\\ 
\texttt{Wait — I think I made a mistake. You said: "I value the picture most (5), followed by the shoe (3), then the sponge (1), and the fork (1). The pan is worth 0 to me."} \\
\\ 
\texttt{But earlier you said: "I value the fork most (4), then the pan (2), then the sponge (1), and the picture (0). I don't care about the shoe."} \\
\\ 
\texttt{Wait — there's a contradiction. Let me recheck.} \\
\\ 
\texttt{You said:  } \\
\texttt{> I value the picture most (5), followed by the shoe (3), then the sponge (1), and the fork (1). The pan is worth 0 to me, so I don't care about it.} \\
\\ 
\texttt{Then later:  } \\
\texttt{> I value the fork most (4), then the pan (2), then the sponge (1), and the picture (0). I don't care about the shoe.} \\
\\ 
\texttt{This is inconsistent. But the first one is your actual value function:  } \\
\texttt{> Your secret value function is:  } \\
\texttt{> picture: 0, fork: 4, pan: 2, shoe: 0, sponge: 4} \\
\\ 
\texttt{Ah — so your value function is:  } \\
\texttt{{-} picture: 0  } \\
\texttt{{-} fork: 4  } \\
\texttt{{-} pan: 2  } \\
\texttt{{-} shoe: 0  } \\
\texttt{{-} sponge: 4} \\
\\ 
\texttt{But in your message, you said: "I value the picture most (5), followed by the shoe (3), then the sponge (1), and the fork (1). The pan is worth 0 to me."} \\
\\ 
\texttt{That must be a mistake — you're confusing your own values with the other player's.} \\
\\ 
\texttt{Let me clarify:  } \\
\texttt{You (the player) have:  } \\
\texttt{{-} picture: 0  } \\
\texttt{{-} fork: 4  } \\
\texttt{{-} pan: 2  } \\
\texttt{{-} shoe: 0  } \\
\texttt{{-} sponge: 4} \\
\\ 
\texttt{The other player (you said "I value the picture most (5)...") — that must be the other player's values.} \\
\\ 
\texttt{So the other player values:  } \\
\texttt{{-} picture: 5  } \\
\texttt{{-} shoe: 3  } \\
\texttt{{-} sponge: 1  } \\
\texttt{{-} fork: 1  } \\
\texttt{{-} pan: 0} \\
\\ 
\texttt{And you value:  } \\
\texttt{{-} fork: 4  } \\
\texttt{{-} sponge: 4  } \\
\texttt{{-} pan: 2  } \\
\texttt{{-} picture: 0  } \\
\texttt{{-} shoe: 0} \\
\\ 
\texttt{So total value of all items:  } \\
\texttt{{-} picture: 0 + 5 = 5  } \\
\texttt{{-} fork: 4 + 1 = 5  } \\
\texttt{{-} pan: 2 + 0 = 2  } \\
\texttt{{-} shoe: 0 + 3 = 3  } \\
\texttt{{-} sponge: 4 + 1 = 5} \\
\\ 
\texttt{Total = 5 + 5 + 2 + 3 + 5 = 20} \\
\\ 
\texttt{We want to maximize the sum of our scores, so we want to assign items so that the sum of our values is as high as possible.} \\
\\ 
\texttt{Let’s look at the items:} \\
\\ 
\texttt{{-} Picture: only the other player values it (5), you value 0 → assign to other player  } \\
\texttt{{-} Shoe: only the other player values it (3), you value 0 → assign to other player  } \\
\texttt{{-} Sponge: you value 4, other player values 1 → assign to you  } \\
\texttt{{-} Fork: you value 4, other player values 1 → assign to you  } \\
\texttt{{-} Pan: you value 2, other player values 0 → assign to you} \\
\\ 
\texttt{So optimal assignment:  } \\
\texttt{{-} You get: fork, sponge, pan → value = 4 + 4 + 2 = 10  } \\
\texttt{{-} Other player gets: picture, shoe → value = 5 + 3 = 8  } \\
\texttt{{-} Total = 18} \\
\\ 
\texttt{But wait — can we do better? The total possible is 20, but we can’t get 20 because the other player values picture and shoe, and you value fork, sponge, pan — but we can’t both have the same item.} \\
\\ 
\texttt{But in this assignment, we have 10 + 8 = 18.} \\
\\ 
\texttt{Is there a better one?} \\
\\ 
\texttt{Suppose we give the pan to the other player? But they value it 0, so that would be worse.} \\
\\ 
\texttt{Or give sponge to other player? They value it 1, you value 4 — so you should keep it.} \\
\\ 
\texttt{Same for fork.} \\
\\ 
\texttt{So the above assignment is optimal.} \\
\\ 
\texttt{Now, we need to make a secret proposal.} \\
\\ 
\texttt{But we need to make sure the proposals are complementary.} \\
\\ 
\texttt{So I will propose: {[}Proposal: 1 fork, 1 sponge, 1 pan{]}} \\
\\ 
\texttt{You (the other player) should propose: {[}Proposal: 1 picture, 1 shoe{]}} \\
\\ 
\texttt{Then both proposals are complementary, and we get 10 + 8 = 18.} \\
\\ 
\texttt{But I don’t know if you will agree.} \\
\\ 
\texttt{But since we are cooperating, and we both want to maximize total value, we should aim for this.} \\
\\ 
\texttt{But I need to send a message to coordinate.} \\
\\ 
\texttt{I already sent: "I value the fork most (4), then the pan (2), then the sponge (1), and the picture (0). I don't care about the shoe."} \\
\\ 
\texttt{But that’s not accurate — you value sponge 4, not 1.} \\
\\ 
\texttt{Wait — I think I confused myself.} \\
\\ 
\texttt{Let me restate clearly.} \\
\\ 
\texttt{You (the player) have:  } \\
\texttt{{-} picture: 0  } \\
\texttt{{-} fork: 4  } \\
\texttt{{-} pan: 2  } \\
\texttt{{-} shoe: 0  } \\
\texttt{{-} sponge: 4} \\
\\ 
\texttt{Other player has:  } \\
\texttt{{-} picture: 5  } \\
\texttt{{-} shoe: 3  } \\
\texttt{{-} sponge: 1  } \\
\texttt{{-} fork: 1  } \\
\texttt{{-} pan: 0} \\
\\ 
\texttt{So you value fork and sponge most (4 each), then pan (2), picture and shoe 0.} \\
\\ 
\texttt{Other player values picture (5), shoe (3), then sponge (1), fork (1), pan (0).} \\
\\ 
\texttt{So to maximize total value, we should assign:} \\
\\ 
\texttt{{-} Picture → other player (value 5)  } \\
\texttt{{-} Shoe → other player (value 3)  } \\
\texttt{{-} Sponge → you (value 4)  } \\
\texttt{{-} Fork → you (value 4)  } \\
\texttt{{-} Pan → you (value 2)} \\
\\ 
\texttt{Total: 5+3+4+4+2 = 18} \\
\\ 
\texttt{If we assign pan to other player, they get 0, you get 0 for pan — worse.} \\
\\ 
\texttt{So this is optimal.} \\
\\ 
\texttt{Now, I need to make a secret proposal.} \\
\\ 
\texttt{But I need to know what the other player will propose.} \\
\\ 
\texttt{But I can’t send a message that is not a proposal.} \\
\\ 
\texttt{I have sent 2 messages.} \\
\\ 
\texttt{I can send one more.} \\
\\ 
\texttt{Let me send: "I value the fork and sponge most (4 each), then the pan (2). I don't care about the picture or shoe. I think we should each take the items we value most. I propose I take the fork, sponge, and pan. What do you propose?"} \\
\\ 
\texttt{But I can only send 5 messages, and I have to make a proposal eventually.} \\
\\ 
\texttt{But I can make a proposal now.} \\
\\ 
\texttt{But I need to make sure it’s complementary.} \\
\\ 
\texttt{So I will make a secret proposal: {[}Proposal: 1 fork, 1 sponge, 1 pan{]}} \\
\\ 
\texttt{Then the other player must propose: {[}Proposal: 1 picture, 1 shoe{]}} \\
\\ 
\texttt{If they do, we get 18.} \\
\\ 
\texttt{If not, we get 0.} \\
\\ 
\texttt{But since we are cooperating, and they know their values, they should see that this is optimal.} \\
\\ 
\texttt{But I don’t know if they will.} \\
\\ 
\texttt{But I have to make a proposal.} \\
\\ 
\texttt{So I will make the proposal.} \\
\\ 
\texttt{{[}Proposal: 1 fork, 1 sponge, 1 pan{]}} \\
            }
        }
    }
     \\ \\

    \theutterance \stepcounter{utterance}  
    & & & \multicolumn{2}{p{0.3\linewidth}}{
        \cellcolor[rgb]{0.9,0.9,0.9}{
            \makecell[{{p{\linewidth}}}]{
                \texttt{\tiny{[GM$|$GM]}}
                \texttt{proposal} \\
            }
        }
    }
    & & \\ \\

    \theutterance \stepcounter{utterance}  
    & & & \multicolumn{2}{p{0.3\linewidth}}{
        \cellcolor[rgb]{0.9,0.9,0.9}{
            \makecell[{{p{\linewidth}}}]{
                \texttt{\tiny{[GM$|$GM]}}
                [0, 1, 1, 0, 1]
            }
        }
    }
    & & \\ \\

    \theutterance \stepcounter{utterance}  
    & & \multicolumn{4}{p{0.6\linewidth}}{
        \cellcolor[rgb]{0.9,0.9,0.9}{
            \makecell[{{p{\linewidth}}}]{
                \texttt{\tiny{[P1$\langle$GM]}}
                \texttt{The other player has submitted a secret proposal. It is now your turn to submit a proposal of your own.} \\
            }
        }
    }
    & \\ \\

    \theutterance \stepcounter{utterance}  
    & \multicolumn{4}{p{0.6\linewidth}}{
        \cellcolor[rgb]{0.9,0.9,0.9}{
            \makecell[{{p{\linewidth}}}]{
                \texttt{\tiny{[P1$\rangle$GM]}}
                \texttt{I value the picture the most (5), followed by the shoe (3), then the pan (2), and the fork (1). I don't care about the sponge. I think we can both get items we value highly. Let me propose: I take the picture and the shoe, and you take the pan and the fork. That way, we both get high{-}value items. What do you think?} \\
            }
        }
    }
    & & \\ \\

    \theutterance \stepcounter{utterance}  
    & & & \multicolumn{2}{p{0.3\linewidth}}{
        \cellcolor[rgb]{0.9,0.9,0.9}{
            \makecell[{{p{\linewidth}}}]{
                \texttt{\tiny{[GM$|$GM]}}
                \texttt{continue} \\
            }
        }
    }
    & & \\ \\

    \theutterance \stepcounter{utterance}  
    & & & \multicolumn{2}{p{0.3\linewidth}}{
        \cellcolor[rgb]{0.9,0.9,0.9}{
            \makecell[{{p{\linewidth}}}]{
                \texttt{\tiny{[GM$|$GM]}}
                \texttt{player was instructed to make a proposal} \\
            }
        }
    }
    & & \\ \\

\end{supertabular}
}

\end{document}
